% Slide — GA: antes (julho)
\begin{frame}{Algoritmo Genético -- \antes}
  \footnotesize Biblioteca \texttt{DEAP}. Representação real $\theta=(\mathrm{ROP}, Q, \mathrm{SS})$.

  \renewcommand{\arraystretch}{1.05}\footnotesize
  \begin{tabular}{@{}p{0.28\linewidth}p{0.66\linewidth}@{}}
    \toprule
    \textbf{Operador} & \textbf{Implementação} \\
    \midrule
    População & 100 indivíduos, inicializados uniformemente nos limites de busca \\
    Gerações & 50 \\
    Seleção & Torneio, tamanho 3 (\texttt{tools.selTournament}) \\
    Cruzamento & \textit{Blend crossover} $\alpha=0{,}5$ (\texttt{tools.cxBlend}), prob. 0,8 \\
    Mutação & Gaussiana, $\sigma=$ maior amplitude do espaço de busca $/10$, prob. 0,05 por gene (\texttt{tools.mutGaussian}) \\
    Elitismo & \textbf{Nenhum} -- \texttt{HallOfFame(1)} só registra o melhor, não o protege do crossover \\
    Avaliação & 1 simulação completa de $T$ ciclos por indivíduo, \textbf{em laço Python} \\
    Aptidão & Soma ponderada $\lambda_1\,\mathrm{NS}-\lambda_2\,\mathrm{CTI}$, pesos fixos \\
    \bottomrule
  \end{tabular}

  \note{
    Este é o inventário completo do GA de julho, operador por operador,
    porque a pergunta natural da banca é "o que exatamente mudou". A
    população, o número de gerações, o tipo de cruzamento e o tipo de
    mutação continuam os mesmos na versão atual -- não houve mudança de
    metaheurística, só de execução e de aptidão. O ponto crítico aqui é a
    ausência de elitismo: sem proteção, o melhor indivíduo encontrado numa
    geração pode ser destruído pelo cruzamento blend na geração seguinte,
    o que explica parte da instabilidade observada entre séries nos
    resultados do Experimento 1. E cada geração custava 100 simulações
    sequenciais em Python puro.
  }
\end{frame}
